\chapter[Sermon 57. Of the Power of the Roman Empire, and of Those Who Worship the Beast: and of the Destruction of Rome, and the Roman Empire.]{Of the Power of the Roman Empire, and of Those Who Worship the Beast: and of the Destruction of Rome, and the Roman Empire.}
\label{sermon:057}
\markright{Sermon 57. Of the Power of the Roman Empire, and of Those Who ...}

And power was given him over all relatives, language, and nation, and all who dwell upon the earth worshipped him, whose names are not written in the Book of Life of the Lamb who was slain from the foundation of the world. If anyone has an ear, let him hear. He who leads into captivity shall go into captivity; he who kills with the sword must be killed with the sword. Here is the perseverance, and the faith of the saints.

\index{Rome}The Apostle, through the revelation of Christ, also speaks of the power and majesty of the Roman Empire. The Roman Empire was indeed of greatest power in the time of Octavian Augustus, also in the time of Domitian, and in the reign of Trajan, also under Hadrian, Aurelian, Diocletian, and Constantine. The greater part of the inhabited world obeyed, namely all of Europe in a manner, Asia, and Africa, as both Latin and Greek histories testify. However, the Lord warns us not to curiously search the counsels of God, being inquisitive, why God gave so great power to the Romans, whom he knew would abuse it to the oppression of Christ’s Church. For where he says that the power was given to Rome, he stills and appeases all murmurings. For empires are of God. But he is most wise, righteous, and holy. Therefore, where he made the kingdoms of the world subject to Rome, he did it wisely, justly, and holily. In that the Romans corrupt God’s ordinance and commit themselves to be governed by the Devil, it comes of evil.

\index{Daniel (Book of)}Let our arguments here cease, for the wise man also says that wicked men and hypocrites rule because of the people’s sins. And that he recounts lineages, languages, and nations, he does in imitation of the Prophet Daniel, who by such phrasing is accustomed to signify a very extensive and powerful empire.

But what does this have to do with us, or what benefit (do you say) comes to us from it, that the Roman Empire is so widely extended throughout the world? Truly, we see how this prophecy has accurately predicted everything that has come before; therefore, there is no room left for doubt about the things that follow. Let us consider moreover that the most powerful kingdoms, which appear invincible to people, may be dissolved by God without any difficulty. Let us therefore learn to fear God, and to walk in his commandments, and to despise these earthly things.

\index{Apocalypse (Book of)}\index{Arethas of Caesarea}\index{Aquinas, Thomas}Now he declares more expressly who will worship the beast: for he said that people in the world will be captivated by the beast and will worship it. He now clarifies this, placing the word “worship” so that it may be understood to apply both to those who are present and to those who are to come. For he speaks not only of people of his time, but of all who, captivated by the admiration of the empire and its majesty, will either deny or despise the faith of Christ. And he says that all who dwell on the earth will worship the beast—and lest anyone refer this absolutely to all, as though none of the true worshippers of God were included, he adds, “whose names are not written in the book of life of the Lamb,” namely, the reprobates, not the elect: the unbelievers. I say, those who despise the word of the gospel, disdain to hear it, and are rebels to Christ. Arethas, the expositor, says they dwell on the earth, those who are unconcerned about heavenly things or the glory that is there, or who devote themselves to earthly pursuits and live a beastly life, according to the same. Thomas Aquinas also brings a testimony from Jeremiah 17: They who depart from me shall be written in the earth. For they have forsaken the fountain of living waters, even the Lord himself. Of the book of life I have spoken in chapters 3 and 5, and will speak of it again in chapters 19 and 20 of the Apocalypse.

Hereunto he adds a noteworthy observation, in the manner of the Apostles, who always, whenever they have occasion to celebrate and explain Christ and the mystery of his redemption, do so. John says that the Lamb has been slain and offered from the beginning of the world. And it is without controversy that by the Lamb is understood Christ.

\index{Hebrews, Epistle to the}It is therefore demanded, how he was slain from the beginning of the world. Many torment themselves at length as they expound, that Christ was slain in Abel, and in all saints, by participation, not by passion. Certainly we may not expound this place after the letter. For Christ could not be slain before he was born. Moreover, the Apostle affirms that Christ, who was from the beginning of the world, has not been slain oftener than once. Read what he says in chapter 9 to the Hebrews. And yet the most true word of God cannot be contrary or repugnant to itself. Therefore, we say, after the common rule of expounding the Scriptures, that the signs have the names of the things signified. For the Lamb was called a Passover, whereof it was a sign. Circumcision was called the league or covenant itself; sacrifices are named sins. So, verily, from the beginning of the world, sacrifices were slain, which were symbols or signs of Christ to be incarnated and offered up once for the cleansing of sins. We understand therefore by this testimony of Christ, that all the sacrifices of the ancient fathers were sacraments of Christ, and that the redemption of Christ has from the beginning of the world been of efficacy to all the faithful. Therefore, this place is notable and worthy to be observed. Hitherto pertains the Apostle’s testimony in 1 Corinthians 10:1: that all our forefathers have eaten of the same spiritual meat with us, and drunk of the same drink, and that they drank of the rock following them, which was Christ.

And up to this point, he has spoken of the majesty of the Roman Empire, its blasphemies, and its sins. Now follows a discussion of the destruction of so great an empire and the punishment of sins. However, this will be addressed again in chapter 17.

And with an acclamation, most commonly used in the gospel, and as it were peculiar to Christ, he stirs up all his listeners, and cries out, “He that has an ear let him hear.” Truly it was to men a wonder, and seemed incredible, that so great a majesty could fall; but yet it falls. The faithful also marveled what should be the end of blasphemies, slaughters, injuries, abominations. Moreover, the doctrine that follows is notable, excellent, and worthy to be kept in memory. Therefore he stirs up all men to attentiveness, and then he says: “Whosoever shall lead into captivity, shall go into captivity; whosoever strikes with the sword, \&c.” For in such sort he declares the destruction of Rome and the Roman empire, that he confirms with all the justice of God’s judgment. And also with a marvelous brevity of God’s sentence given or pronounced against Rome, he partakes of that immeasurable power. And this is both by the law of God, by the law of nature, and by the law of all nations received as a thing most just, that every man should look to have the same done to him that he doth to another. For to this belongs the sentence rehearsed of Noah in Genesis 9: “He that sheds blood, his blood shall be shed.” The same is repeated in Leviticus 33: “Woe to him that spoils, shall not thou be spoiled?” A testimony whereof is Nineveh with the prophet Nahum, and Babylon with all the prophets. Therefore has the Lord taught in the gospel, “Whatsoever you would that men should do to you, do you the same unto them also.” With what measure you mete unto others, with the same shall others mete unto you again. Whosoever strikes with the sword, with the sword shall perish. Therefore it is most reasonable, that since Rome has spoiled the whole world, and injured all nations, and made cruel war upon all men, it should be again of all nations invaded, spoiled, torn, and trodden under foot. Let us mark this judgment of God, and let us fear God, and do good unto men. For here is sentence given against all men that do injury to their neighbors, but especially those which invade innocents with unjust wars, and which they are hired to make, \&c.

\index{Jerome, Saint}And here we must repeat something from Histories, whereby the truth of this prophecy may be better known and understood. When the most excellent Prince Constantine received the government of the empire, as if abhorring Rome, he built Constantinople and made it the seat of the empire. And from that time the majesty of Rome began to fall into ruin. Under the emperor Gratian, a prince most witty, the Barbarians were a great terror to the Romans, whereupon Gratian made a league with them. Stilicho, father-in-law to Honorius, a Vandal by birth, diminished the wages of the Gothians and other league fellows of the people of Rome; for which cause they took armor. Yet being pacified again, they were stirred up afterward through the malice of Stilicho and of Duke Saule, and under the conduct of Athalaric their king, they marched to Rome, laid siege to it, and besieged it for the space of two years, at length taking and spoiling it. This siege and spoil Saint Jerome in his Epistle bewails much. Orosius writes much and Christianly hereof in the 29th chapter of the 7th book of Histories. It is reported that Rome was taken the first day of April, in the year 412. Yet the Gothians immediately leaving the city, removed into other places thereby; nevertheless, being again inflamed with fury, they returned, and under their captain Athaulphus, they plagued and spoiled Rome worse than they did before. The king had determined, extinguishing the name of Romans, to have called the city Gothia, if he had not been dissuaded by Galla Placidia, daughter to Honorius. A few years after, Rome was taken again by Gensericus, king of the Vandals; and that which was enriched and replenished with the robberies of all nations was by fourteen days together emptied clean. After came Odacer with the Germans; and putting down the name of emperor, reigned over the city himself as king for the space of 15 years. Whom Theodoric of Verona expelled and slew. And there reigned with his Ostrogoths about 50 years. Then it was recovered by Belisarius for Justinian, emperor of Greece, but to the utter destruction of Rome. For Totila, king of Gothia, discomfited both the Greek and Roman army at Placentia; after he besieged Rome, scaled, took, sacked, overthrew, and set it on fire. The city burned thirteen days. Neither was there any man in it for the space of forty days. Read the 4th book of Sabellicus the 8th Aeneade. Perhaps I shall discourse more at large of the destruction of Rome in the 17th chapter. Wherefore within the space of 136 years, Rome came seven times into strangers' hands, and was sacked most cruelly, and fell by the edge of the sword, and was led into captivity: which hath long stricken with the sword, and led away all nations prisoners. This was the just judgment of God.

\index{John, Saint}And Saint John adds a doctrine, how the faithful should behave in such great troubles and adversities. That is to say, while the Romans rule and rage, also in those bloody and cruel changes and destruction of the Roman Empire, the saints will need patience, or perseverance, and faith. These two virtues will preserve the faithful, so that they do not perish as well. Concerning patience, the Lord speaks in Saint Luke, chapter 21: “In your patience you will possess your souls.” Blessed John speaks of faith: “And this is the victory that overcomes the world, even your faith.” Impatience and incredulity have led many into the denying of the faith, to idolatry and to all ungodliness. Thus, we also learn how to arm ourselves in our days against all ungodliness. The Lord deliver us from evil. Amen.
